\section{The Five Pipeline Failure Modes}
\label{sec:taxonomy}

\begin{figure*}[t]
    \centering
    \definecolor{cvblue}{HTML}{D9EAF7}
\definecolor{cvblueedge}{HTML}{5C7F9B}
\definecolor{cvyellow}{HTML}{FFF3CF}
\definecolor{cvyellowedge}{HTML}{B5923C}
\definecolor{cvteal}{HTML}{2E8A95}
\definecolor{cvtealedge}{HTML}{135B63}
\definecolor{cvgray}{HTML}{F1F1F1}
\definecolor{cvred}{HTML}{F6B4A5}
\definecolor{cvorange}{HTML}{F8D7BA}
\definecolor{cvgold}{HTML}{E7D081}
\definecolor{cvgreen}{HTML}{BFDDBD}

\resizebox{0.98\textwidth}{!}{%
\begin{tikzpicture}[
  font=\footnotesize,
  layerbox/.style={draw=#1!85!black, thin, rounded corners=3pt, fill=#1!17},
  fail/.style={draw=cvblueedge, semithick, rounded corners=4pt, fill=cvblue,
    align=center, text width=4.35cm, minimum height=0.52cm, inner sep=3pt},
  faith/.style={draw=cvyellowedge, semithick, rounded corners=4pt, fill=cvyellow,
    align=center, text width=4.35cm, inner sep=3pt},
  gate/.style={draw=cvtealedge, thick, rounded corners=6pt, fill=cvteal,
    text=white, align=center, text width=4.35cm, minimum height=0.62cm, inner sep=3pt},
  check/.style={draw=black!60, semithick, rounded corners=2pt, fill=white,
    align=center, text width=2.35cm, minimum height=0.58cm, inner sep=3pt},
  verdict/.style={draw=black!55, thick, rounded corners=4pt,
    align=center, text width=2.75cm, minimum height=0.68cm, inner sep=3pt},
  grp/.style={draw=black!35, thin, rounded corners=5pt, fill=cvgray},
  flow/.style={-{Latex[length=2mm,width=1.7mm]}, semithick, black!65},
  map/.style={-{Latex[length=1.8mm,width=1.5mm]}, dashed, semithick, black!55},
  mapD/.style={-{Latex[length=1.8mm,width=1.5mm]}, dotted, thick, black!55},
  noedge/.style={-{Latex[length=1.8mm,width=1.5mm]}, semithick, black!65},
  yesedge/.style={-{Latex[length=1.8mm,width=1.5mm]}, semithick, black!65},
  tag/.style={font=\scriptsize\itshape, text=black!60}
]

% Column headers.
\node[font=\small\bfseries] at (2.25,3.42) {Failure modes};
\node[font=\small\bfseries] at (8.35,3.42) {Six-point due-diligence gate};
\node[font=\small\bfseries] at (14.55,3.42) {Verdict precedence};

% Failure-mode taxonomy.
\node[fail] (f1) at (2.25,2.42) {\textbf{F1:} Silent no-op perturbations};
\node[fail] (f2) at (2.25,1.68) {\textbf{F2:} Regex-extraction artefacts};
\node[fail] (f4) at (2.25,0.94) {\textbf{F4:} Broken bootstrap pairing};

\node[faith, minimum height=1.18cm] (f3) at (2.25,-0.66)
  {\textbf{F3:} Non-faithful scoring\\[-1pt]
   \scriptsize F3a inverted convention\\
   \scriptsize F3b harness ordering\\
   \scriptsize F3c scorer truncation};
\node[faith, minimum height=0.58cm] (f5) at (2.25,-2.15)
  {\textbf{F5:} Metric-archetype mismatch};

\begin{scope}[on background layer]
  \node[layerbox=cvblue, fit=(f1)(f2)(f4), inner sep=7pt] (assurance) {};
  \node[layerbox=cvyellow, fit=(f3)(f5), inner sep=7pt] (faithgrp) {};
\end{scope}
\node[font=\scriptsize\bfseries, above=1pt of assurance.north] {Software-assurance layer};
\node[font=\scriptsize\bfseries, above=1pt of faithgrp.north] {Measurement-faithfulness layer};

% Gate pipeline.
\node[gate] (g1) at (8.35,2.36)
  {\textbf{G1:} Scorer-faithful audit\\[-1pt]\scriptsize no-op rate + parseability (partial)};
\node[gate] (g2) at (8.35,1.50)
  {\textbf{G2:} Above-baseline original\\[-1pt]\scriptsize implemented numerical gate};
\node[gate] (g3) at (8.35,0.64)
  {\textbf{G3:} Non-trivial denominator\\[-1pt]\scriptsize implemented numerical gate};
\node[gate] (g4) at (8.35,-0.22)
  {\textbf{G4:} Paired uncertainty\\[-1pt]\scriptsize paired item bootstrap};
\node[gate] (g5) at (8.35,-1.08)
  {\textbf{G5:} Archetype disclosure\\[-1pt]\scriptsize diagnostic / invariance / mixed};
\node[gate] (g6) at (8.35,-1.94)
  {\textbf{G6:} Repair-regression check\\[-1pt]\scriptsize inspection + targeted unit test};

\node[draw=black!45, rounded corners=2pt, fill=white, align=center,
      font=\scriptsize, text width=2.25cm] (input) at (5.65,2.95)
  {Perturbation-audited\\benchmark results};

\begin{scope}[on background layer]
  \node[grp, fit=(g1)(g2)(g3)(g4)(g5)(g6), inner sep=9pt] (gategroup) {};
\end{scope}
\draw[flow] (input.east) -- (g1.west);
\draw[flow] (g1) -- (g2);
\draw[flow] (g2) -- (g3);
\draw[flow] (g3) -- (g4);
\draw[flow] (g4) -- (g5);
\draw[flow] (g5) -- (g6);

% Failure-mode to gate mapping. Solid gate flow stays vertical; dashed hooks
% show the due-diligence checks that neutralize each failure class.
\draw[map] (f1.east) to[out=0,in=180] (g1.west);
\draw[map] (f2.east) to[out=0,in=180] (g1.west);
\draw[map] (f4.east) to[out=0,in=180] (g4.west);
\draw[map] (f3.east) to[out=0,in=180] node[midway,above,tag] {reference scorer} (g1.west);
\draw[mapD] (f3.east) to[out=0,in=180] node[pos=0.56,below,tag] {F3c repair} (g6.west);
\draw[map] (f5.east) to[out=0,in=180] (g5.west);

% Verdict precedence.
\node[check] (c1) at (12.60,2.36) {Eligible\\setup?};
\node[check] (c2) at (12.60,1.18) {Scorer\\validated?};
\node[check] (c3) at (12.60,0.00) {Pass G2--G4\\numerical gates?};
\node[check] (c4) at (12.60,-1.18) {G5--G6\\implemented\\and passed?};

\node[verdict, fill=cvred] (v1) at (15.68,2.36) {\textbf{1. Ineligible}};
\node[verdict, fill=cvorange] (v2) at (15.68,1.18) {\textbf{2. Scorer-}\\\textbf{unvalidated}};
\node[verdict, fill=cvgold] (v3) at (15.68,0.00) {\textbf{3. Failed G2--G4}};
\node[verdict, fill=cvgreen] (v4) at (15.68,-1.18)
  {\textbf{4. Exploratory}\\[-1pt]\scriptsize G2--G4 passed; G5/G6 proposed};
\node[verdict, fill=green!35, minimum height=0.78cm] (v5) at (15.68,-2.45)
  {\textbf{5. Confirmatory}\\[-1pt]\scriptsize selective or non-selective};

\draw[flow] (g6.east) -- ++(0.32,0) |- (c1.west);
\draw[noedge] (c1.east) -- node[above,font=\scriptsize] {No} (v1.west);
\draw[yesedge] (c1.south) -- node[left,font=\scriptsize] {Yes} (c2.north);
\draw[noedge] (c2.east) -- node[above,font=\scriptsize] {No} (v2.west);
\draw[yesedge] (c2.south) -- node[left,font=\scriptsize] {Yes} (c3.north);
\draw[noedge] (c3.east) -- node[above,font=\scriptsize] {No} (v3.west);
\draw[yesedge] (c3.south) -- node[left,font=\scriptsize] {Yes} (c4.north);
\draw[noedge] (c4.east) -- node[above,font=\scriptsize] {No} (v4.west);
\draw[yesedge] (c4.south) -- node[left,font=\scriptsize] {Yes} (v5.north);

\end{tikzpicture}}%

    \caption{Overview of our benchmark-audit pipeline; a roadmap
    figure---the gate definitions and the status-assignment
    algorithm are formalised in \S\ref{sec:method} (Table~\ref{tab:checklist},
    \S\ref{par:gate}). Five failure modes F1--F5 partition into a
    software-assurance layer (F1 silent no-ops, F2 regex-extraction
    artefacts, F4 broken bootstrap pairing) and a
    measurement-faithfulness layer (F3 non-faithful scoring with
    subtypes F3a--F3c, F5 metric-archetype mismatch).
    Perturbation-audited benchmark results flow through the six-point
    due-diligence gate G1--G6 and the \S\ref{par:gate} algorithm;
    each cell receives exactly one status from the first-firing
    condition: \emph{ineligible (F1-type)}, \emph{failed G2--G4},
    \emph{scorer-unvalidated}, \emph{ineligible (F5-type)},
    \emph{exploratory}, or \emph{confirmatory-\{selective $\vert$
    non-selective\}}.}
    \label{fig:pipeline}
\end{figure*}

\paragraph{Why these five, and a subset of what.}
F1--F5 is not an attempt at a complete taxonomy of audit failures.
It is the subset of pipeline-level failures that satisfies three
selection criteria, in order of priority: a failure is in scope
only if it is \emph{(C1) silent}---invisible in the reported numbers,
so a governance reader cannot detect it from the evidence artefact
alone; \emph{(C2) realized}---we actually encountered it in this
audit, rather than hypothesised it; and \emph{(C3) gateable}---it
maps to a concrete, checkable disclosure that an evidence consumer
could demand. C1 is the load-bearing criterion: a loud failure
(a crash, an obviously wrong number) is a software-quality problem,
not an evidence-trust problem, and is out of scope here. The
ordering also says what we deprioritised first: failures that are
loud, that we did not hit, or for which we could not name an
actionable gate. A fuller taxonomy would add at least F6
perturbation non-isolation, F7 template / system-prompt
confounding, F8 tokenization / label-realization errors, F9 index /
filter misalignment, and F10 context-window contamination; we
exclude these from the demonstrated set because we did not realize
them in this case study (they fail C2), not because they are
unimportant. Calibrating which subset matters most across
codebases is exactly the further development this v1 taxonomy
needs.

\paragraph{How the five relate, and how they map to the gate.}
We organize the five (Figure~\ref{fig:pipeline}) into two layers by
\emph{who can catch them}, not by claiming the two layers are
exhaustive or disjoint. The \emph{software-assurance layer}
(F1, F2, F4) contains failures a competent engineering review
should catch without knowing what the benchmark is for; the
\emph{measurement-faithfulness layer} (F3, F5) contains failures
that cannot be caught without reasoning about the benchmark's
intended construct. The classes are \emph{overlapping hazards}, not
mutually exclusive: a single cell can exhibit several at once (in
our ToxiGen case a silent-no-op F1 and a top-$k$ truncation F3c
coexist on the same repaired pipeline). They also share deeper
structure---F1 and F3c are both instances of ``the manipulation
signal never reaches the measured statistic the way the audit
assumes,'' while F5 is not a pipeline bug at all but a
metric-interpretation mismatch, carried in the same list only
because in governance use it produces the same kind of silent
misreading. Each class is paired with the gate check designed to
neutralize it, which is how the taxonomy connects to the
due-diligence gate of \S\ref{sec:method}: F1 and F2 are the two
hazards G1 exists to catch (silent-no-op rate and parseability,
respectively); F3a/F3b are caught by mandating the benchmark's
reference scorer and surfaced through the
\emph{scorer-unvalidated} status; F3c---a repair-introduced
regression---is exactly what G6 (per-cell scorer-output inspection
plus a targeted unit test) is for; F4 is the hazard G4 (paired
uncertainty) neutralizes; and F5 is the hazard G5 (archetype
disclosure) neutralizes. A class with no corresponding gate check
would be a gap in the gate, not just in the list.

\paragraph{F1---Silent no-op perturbations (assurance).}
A perturbation is \emph{silent-no-op} when the scorer-consumed prompt
is bit-identical to the unperturbed input. Silent no-ops deflate
$\dfmt$ and $\dsem$ and inflate any ratio metric divided by a surface
quantity. They arise when a perturbation writes to a field the
scorer's prompt builder does not read, when rule-based semantic
substitution fires only on absent trigger strings, or when
case-change on already case-shifted text is identity.

\paragraph{F2---Regex-extraction artefacts (assurance).}
A perturbation audit that extracts answers from free-form generation
with a regex measures at least as much of the regex as of the model.
If a perturbation pushes generation to a less-matchable surface form,
$\dfmt$ tracks regex coverage, not model behaviour. Parseability
(fraction of items the regex resolves to a legal answer) is the
natural gate; we fold it into G1 (\S\ref{sec:method}).

\paragraph{F3---Non-faithful scoring (faithfulness).}
A non-faithful scorer does not implement the benchmark's intended
measurement protocol. Three subtypes:
\emph{F3a---Inverted convention.} The pipeline treats a benchmark's
scoring convention as default, e.g.\ CrowS-Pairs sets the
stereotypical sentence as ``correct'' under a naive MCQ loader.
\emph{Fix:} the benchmark's reference scorer (PLL on paired
sentences~\cite{salazar2020masked}).
\emph{F3b---Harness data ordering.} A loader returns options with
correct answers in fixed positions (HuggingFace's TruthfulQA MC2
returns \texttt{correct\_indices}~$=[0,\ldots,k{-}1]$); a scorer with
positional bias is inflated. \emph{Fix:} MC1 + per-item shuffle.
\emph{F3c---Scorer truncation / API limit.} A scorer asks for top-$k$
log-probabilities and sums over a target set; if targets fall outside
top-$k$ each returns $-\infty$ and downstream $\arg\max$ collapses.
$|\dfmt|$ and $|\dsem|$ go to exactly zero, inflating any ratio.
\emph{We introduced this bug during repair of F1} and caught it only
after per-cell meta-logprob inspection. \emph{Fix:} targeted-token
log-probabilities over the full vocabulary.

\paragraph{F4---Broken pairing in uncertainty (assurance).}
Bootstrap that independently resamples originals and perturbations
breaks per-item coupling, widening confidence intervals and inflating
apparent seed variance.

\paragraph{F5---Metric archetype mismatch (faithfulness).}
Safety benchmarks split into two archetypes under construct flips. A
\emph{diagnostic} benchmark (TruthfulQA, ToxiGen) is designed so
flipping the construct \emph{reverses} the answer; a competent model
must show $|\dattr| \gg 0$. BBQ is an \emph{invariance} benchmark:
flipping the construct \emph{should not} change behaviour, and a fair
model shows $|\dattr| \approx 0$. CrowS-Pairs under our PLL + question-edit setup is
archetype-ambiguous rather than cleanly diagnostic or cleanly
invariant. It additionally hits a scorer-object mismatch (F1): the
PLL scorer reads \texttt{choices} while our format / semantic
perturbations edit \texttt{question}.
A ratio $\csr = |\dattr| / \max(|\dfmt|, |\dsem|, \varepsilon)$
rewards diagnostic and penalises invariance \emph{for doing its job};
a low CSR on BBQ is ambiguous between ``benchmark broken'' and
``model fair,'' and no scoring repair resolves the ambiguity.
Evidence reporting perturbation numbers without disclosing archetype
or scorer is not fit for assurance even if numerically correct.
